% Domain Kit Manual Template
% Replace Home Domain Kit with the human-readable kit name (e.g., "DRI Domain Kit")
% Replace home with the URL-safe slug (e.g., "dri")
% Run: make init DOMAIN_NAME="DRI Domain Kit" DOMAIN_SLUG=dri

\documentclass[
	11pt,
	oneside,
	english,
	singlespacing,
	headsepline,
]{MastersThesis}

\usepackage[utf8]{inputenc}
\usepackage[T1]{fontenc}
\usepackage{microtype}
\usepackage{hyperref}
\usepackage{enumitem}
\usepackage{longtable}
\usepackage{booktabs}
\usepackage{graphicx}
\usepackage{xcolor}
\usepackage{amssymb}
\usepackage{tcolorbox}
\usepackage{newunicodechar}

% ── Diagrams (TikZ) ──────────────────────────────────────────────────────────
\usepackage{tikz}
\usetikzlibrary{arrows.meta, automata, positioning, shapes.geometric,
                fit, chains, calc, backgrounds}

% Node styles
\tikzset{
  stepbox/.style={
    rectangle, rounded corners=3pt,
    draw=black!60, fill=white,
    font=\small\ttfamily,
    minimum width=2.2cm, minimum height=0.7cm,
    align=center, inner sep=4pt
  },
  catbox/.style={
    rectangle, rounded corners=3pt,
    draw=black!80, fill=black!8,
    font=\small\bfseries,
    minimum width=2.4cm, minimum height=0.75cm,
    align=center, inner sep=4pt
  },
  gertarrow/.style={
    ->, >=Stealth, thick, draw=black!60
  },
  gertfit/.style={
    draw=black!30, dashed, rounded corners=4pt,
    inner sep=6pt
  }
}

% ── Syntax Highlighting (minted) ─────────────────────────────────────────────
\usepackage{minted}
\setminted{
  fontsize=\small,
  baselinestretch=1.1,
  breaklines=true,
  breakanywhere=false,
  autogobble=true,
}
\setminted[yaml]{
  style=friendly,
  linenos=false,
  frame=leftline,
  framesep=6pt,
  rulecolor=\color{black!25}
}
\setminted[go]{
  style=friendly,
  linenos=false,
  frame=leftline,
  framesep=6pt,
  rulecolor=\color{black!25}
}
\setminted[json]{
  style=friendly,
  linenos=false,
  frame=leftline,
  framesep=6pt,
  rulecolor=\color{black!25}
}

% Convenience command for inline code
\newcommand{\yinline}[1]{\mintinline{yaml}{#1}}
% Placeholder macro: renders unfilled template placeholders visibly during development.
% Replace \ph{NAME} with actual content when filling in this domain kit manual.
\newcommand{\ph}[1]{\textcolor{orange!80!black}{\texttt{<#1>}}}

% Define Unicode characters for checkmarks and tree symbols
\newunicodechar{✓}{\checkmark}
\newunicodechar{✗}{$\times$}
\newunicodechar{└}{L}
\newunicodechar{─}{-}
\newunicodechar{├}{+}
\newunicodechar{│}{|}

\hypersetup{
  colorlinks=true,
  linkcolor=blue,
  urlcolor=blue,
  citecolor=blue,
  pdftitle={Home Domain Kit Manual},
  pdfauthor={The gert Project}
}

\thesistitle{Home Domain Kit Manual}
\author{The gert Project}
\supervisor{Technical Documentation}
\degree{User Manual}
\university{Gert Project}
\department{Domain Kits}
\subject{Home Domain Kit}
\keywords{home, gert, domain kit, yaml, dsl}

\begin{document}

\frontmatter

\maketitle

\tableofcontents
\newpage

\mainmatter

\chapter{Introduction}
\label{ch:introduction}

% ─────────────────────────────────────────────────────────────────────────────
\section{What This Manual Covers}
\label{sec:intro-covers}
% ─────────────────────────────────────────────────────────────────────────────

This manual is the authoritative reference for the \texttt{gert.home} Domain Kit---a
specialized vocabulary layer for authoring household maintenance runbooks within the gert
framework. The kit implements the \textbf{Property} model, providing step types, scheduling
primitives, and evidence-capture patterns designed for homeowners and property managers tracking
household maintenance workflows.

The manual covers:

\begin{itemize}
  \item The \texttt{Property} model and role taxonomy
  \item \texttt{gert.home} step types: \texttt{home.routine},
        \texttt{home.incident}, \texttt{home.delegate}
  \item Delegation and away-mode enforcement via \texttt{IsAwayModeActive()}
  \item Evidence types: \texttt{photo}, \texttt{note}, \texttt{checklist}, \texttt{none}
  \item Consumable tracking: replacement schedules and stock alerts
  \item Complete field reference for all \texttt{gert.home} schema elements
\end{itemize}

This is \emph{not} a general introduction to gert or Domain Kits. For core gert concepts
(runbook structure, core step types, the Domain Kit model), see the \emph{gert Design
Document}. For guidance on developing your own Domain Kit, see the \emph{Domain Kit Development
Guide}.

% ─────────────────────────────────────────────────────────────────────────────
\section{Audience and Prerequisites}
\label{sec:intro-audience}
% ─────────────────────────────────────────────────────────────────────────────

This manual is written for:

\begin{description}
  \item[Homeowners] --- The primary audience. Homeowners manage their own property file,
        define zones and assets, schedule routines, and review evidence after each completed
        task.
  \item[Property managers] --- Professionals managing multiple properties. They author one
        \texttt{.home.yaml} file per property, configure delegations for caretakers, and
        review compliance bundles.
  \item[Household contractors] --- Gardeners, plumbers, and other tradespeople who are
        assigned routines via \texttt{executor.role} hints or active delegation blocks. They
        execute tasks and submit evidence but do not modify the property file.
\end{description}

\paragraph{Prerequisites.}
You should be familiar with:

\begin{itemize}
  \item \textbf{gert core concepts} --- runbook structure, core step types (\texttt{cli},
        \texttt{manual}, \texttt{tool}, \texttt{branch}, \texttt{iterate}), variables, and the
        execution model. See \emph{gert Design Document}, §1--§3.
  \item \textbf{Domain Kit concepts} --- what a Domain Kit is, how Kit-specific step types
        compile (lower) into core primitives, and the relationship between authoring vocabulary
        and runtime execution. See \emph{gert Design Document}, §4.
  \item \textbf{YAML syntax} --- \texttt{gert.home} property files are authored in YAML. You
        do not need to know JSON Schema or Go.
\end{itemize}

% ─────────────────────────────────────────────────────────────────────────────
\section{What is a Property?}
\label{sec:intro-concept}
% ─────────────────────────────────────────────────────────────────────────────

\textbf{Property} is the top-level household unit in \texttt{gert.home}. It is the single
source of truth for everything that must be maintained: the named areas (Zones), tracked
physical items (Assets), recurring tasks (Routines), reactive repair workflows
(IncidentTemplates), items on replacement schedules (Consumables), and temporary assignees
(Delegations).

\paragraph{Origin.}
The Property model is drawn from household management systems and facilities management
software. The core insight is that maintenance accountability requires an explicit owner and a
structured record---not just a to-do list. gert.home brings that structure into the runbook
execution engine.

\paragraph{Why the Property Model Matters.}
Without the Property model, household maintenance workflows face:
\begin{itemize}
  \item \textbf{Reactive chaos} --- tasks are remembered only when something breaks, not on
        a schedule.
  \item \textbf{No audit trail} --- there is no record of who did what, when, or what
        evidence was captured.
  \item \textbf{Missed delegations} --- when the homeowner is away, there is no explicit
        assignment of responsibility to a delegate.
\end{itemize}

The Property model resolves this by making household accountability explicit and
machine-readable through the runbook schema.

\paragraph{The Property in \texttt{gert.home}.}
The \texttt{gert.home} Domain Kit implements the Property model through:
\begin{itemize}
  \item A structured \texttt{property} block: \texttt{name}, \texttt{zones[]}
  \item Tracked items: \texttt{assets[]}, \texttt{consumables[]}
  \item Scheduled work: \texttt{routines[]} with cadence and evidence requirements
  \item Reactive repair: \texttt{incident\_templates[]} with sequential and parallel steps
  \item Away-mode: \texttt{delegations[]} with scoped assignments and permissions
\end{itemize}

% ─────────────────────────────────────────────────────────────────────────────
\section{How \texttt{gert.home} Implements the Property Model}
\label{sec:intro-implementation}
% ─────────────────────────────────────────────────────────────────────────────

The \texttt{gert.home} kit extends gert core with domain-specific step types and
metadata fields. These compile down to core primitives at build time so the gert runtime
never sees \texttt{gert.home}-specific vocabulary.

\subsection{Step Types}

\begin{description}
  \item[\texttt{home.routine}] --- A scheduled recurring maintenance task. Lowers to an
        \texttt{iterate} step containing a \texttt{manual} step. Supports fixed-interval
        (\texttt{every: 7d}) and seasonal cadences.
  \item[\texttt{home.incident}] --- A reactive repair workflow with sequential, parallel, and
        branching steps. Each incident step lowers to a \texttt{manual}, \texttt{branch}, or
        \texttt{parallel} core step as appropriate.
  \item[\texttt{home.delegate}] --- An away-mode delegation block. Checked at run-time via
        \texttt{IsAwayModeActive()}; when active, substitutes the delegate identity as executor
        for the assigned routines or zones.
\end{description}

\subsection{Metadata Fields}

\texttt{gert.home} property files include additional top-level blocks:

\begin{description}
  \item[\texttt{property}] --- Display name and zone list for the household unit.
  \item[\texttt{assets}] --- Tracked physical items with installation date and warranty info.
  \item[\texttt{consumables}] --- Items on a replacement schedule with stock tracking.
  \item[\texttt{delegations}] --- Away-mode assignment blocks with date ranges and permissions.
\end{description}

\subsection{Evidence Types}

\texttt{gert.home} defines four evidence types for routine and incident steps:

\begin{description}
  \item[\texttt{photo}] --- A photograph required to prove task completion (e.g., clear pool
        water, mowed lawn).
  \item[\texttt{note}] --- A free-text note from the executor (e.g., cost of repair,
        measurements recorded).
  \item[\texttt{checklist}] --- A structured checklist that must be completed item by item
        before the step closes.
  \item[\texttt{none}] --- No evidence required; the step is self-certifying.
\end{description}

\subsection{Delegation Integration}

When a \texttt{home.delegate} block is active, the \texttt{gert.home} runtime calls
\texttt{IsAwayModeActive()} before each routine execution. If the delegation is active and the
routine or zone is in the delegate's \texttt{assigns} list, the delegate's identity is
substituted as executor. When the delegation expires, the homeowner resumes accountability
automatically.

% ─────────────────────────────────────────────────────────────────────────────
\section{Document Structure}
\label{sec:intro-structure}
% ─────────────────────────────────────────────────────────────────────────────

The remaining chapters are organized as follows:

\begin{description}
  \item[Chapter 1: gert.home Concepts] --- The Property model, role taxonomy, the
        delegation chain, and cadence/SLA concepts.
  \item[Chapter 2: Schema Reference] --- The \texttt{gert.home} kit schema: step types,
        property metadata fields, and evidence types.
  \item[Chapter 3: Authoring Guide] --- How to write a \texttt{.home.yaml} property file
        from scratch, with a complete pool-maintenance example.
  \item[Chapter 4: gert Integration] --- Role-based gating, away-mode enforcement, cadence
        checking, and the audit trail.
  \item[Chapter 5: Evidence and Tracing] --- The evidence model, run directory structure,
        timeline projection, SHA-256 verification, and compliance bundles.
  \item[Chapter 6: Reference] --- Complete field reference tables for all
        \texttt{gert.home} schema elements and environment variables.
\end{description}

\chapter{gert.home Concepts}
\label{ch:concepts}

% ─────────────────────────────────────────────────────────────────────────────
\section{The Property Model}
\label{sec:concepts-model}
% ─────────────────────────────────────────────────────────────────────────────

The \textbf{Property} model is the organizing principle of \texttt{gert.home}: a single
household unit, owned by one accountable homeowner, that holds every zone, asset, routine,
and incident template as structured data. At any moment during a runbook execution, the
Property model makes it clear what needs doing, who is responsible, and what evidence must
be captured.

\paragraph{Single-Owner Invariant.}
A Property has exactly one \texttt{homeowner} at any given time. The homeowner is accountable
for all routines and incidents by default. Ownership does not transfer during normal operation;
only an active \texttt{delegations} block substitutes a delegate as executor for a scoped
set of tasks while the homeowner is away.

\paragraph{Property vs.\ Executor.}
The Property defines \emph{what needs doing} --- zones, assets, routines, and incidents.
\texttt{gert.home} defines \emph{how it is done} --- step types, evidence requirements, and
delegation rules. The homeowner is accountable for the Property's state; the executor (which
may be a delegate or a contracted tradesperson) is responsible for completing individual steps.

% ─────────────────────────────────────────────────────────────────────────────
\section{Role Taxonomy}
\label{sec:concepts-roles}
% ─────────────────────────────────────────────────────────────────────────────

The \texttt{gert.home} kit defines two governance roles and one advisory executor hint.

\subsection{homeowner}

The \texttt{homeowner} is the primary responsible party for the property. They are accountable
for all routines and incidents by default and are the default executor when no delegation is
active. The homeowner authors the property file and configures delegation blocks.

\begin{minted}{yaml}
# Homeowner identity is implicit — it is the authenticated gert user
# running the property file. No explicit declaration is required.
property:
  name: Casa Santiago
\end{minted}

\begin{description}
  \item[\textbf{Accountability}] The homeowner is accountable for every routine and incident
        in the property file, even when a delegate is active.
  \item[\textbf{Default executor}] When no delegation is active, all routines are assigned
        to the homeowner.
  \item[\textbf{Override}] The homeowner can always execute any step, regardless of
        \texttt{executor.role} hints.
\end{description}

\subsection{delegate}

A \texttt{delegate} is a temporary assignee activated by a \texttt{delegations} block with a
defined date range. The delegate's scope is limited to the routines and zones listed in the
\texttt{assigns} block. Permissions are explicitly declared per delegation.

\begin{minted}{yaml}
delegations:
  - delegate:
      name: Carlos Méndez
      contact: "+56 9 8765 4321"
    active:
      from: 2025-04-25
      to: 2025-05-02
    assigns:
      - routine: water_plants
      - zone: pool
    permissions:
      can_report_incidents: true
      can_modify_routines: false
      can_view_history: true
\end{minted}

\begin{description}
  \item[\texttt{assigns}] Scopes the delegation to specific routines or zones. Tasks outside
        this list remain with the homeowner.
  \item[\texttt{permissions}] Declares what the delegate may do beyond executing assigned
        tasks.
  \item[\textbf{Automatic expiry}] When \texttt{active.to} passes, the delegation expires
        and all accountability reverts to the homeowner automatically.
\end{description}

\subsection{Executor Hints}

\texttt{executor.role} is an advisory annotation on a \texttt{home.routine} step. It signals
which type of tradesperson should perform the task (e.g., \texttt{gardener},
\texttt{plumber}, \texttt{contractor}). This is informational only --- the runtime does not
block execution based on executor role in \texttt{home/v0}.

\begin{tcolorbox}[colback=yellow!10,colframe=orange!70,title=Note]
  \texttt{home/v0} does not enforce executor role hints at runtime. Role gating is a planned
  feature for \texttt{home/v1}. The \texttt{executor.role} field is advisory: it is shown in
  the task UI and included in notification messages, but does not prevent the homeowner or
  delegate from executing any step.
\end{tcolorbox}

\begin{minted}{yaml}
routines:
  - id: lawn_mow
    label: Mow front lawn
    zone: front_lawn
    cadence:
      seasonal:
        summer: 7d
        autumn: 10d
        winter: 21d
        spring: 10d
    executor:
      role: gardener
\end{minted}

% ─────────────────────────────────────────────────────────────────────────────
\section{Role Assignment in Property Files}
\label{sec:concepts-role-assignment}
% ─────────────────────────────────────────────────────────────────────────────

Roles are expressed in two locations in the property file.

\paragraph{Runbook-level delegations block.}
The \texttt{delegations} list at the top level declares all away-mode assignments. Each entry
has a date range, an assigns scope, and a permissions map:

\begin{minted}{yaml}
delegations:
  - delegate:
      name: Carlos Méndez
      contact: "+56 9 8765 4321"
      email: carlos@example.com
    active:
      from: 2025-04-25
      to: 2025-05-02
    assigns:
      - routine: pool_clean
      - zone: front_lawn
    permissions:
      can_report_incidents: true
      can_modify_routines: false
      can_view_history: true
    notifications:
      remind_delegate_hours_before: 24
      notify_owner_on_completion: true
\end{minted}

\paragraph{Routine-level executor hint.}
Individual routines may carry an \texttt{executor.role} advisory annotation:

\begin{minted}{yaml}
routines:
  - id: pool_clean
    label: Pool cleaning
    zone: pool
    cadence:
      every: 7d
    executor:
      role: gardener
    evidence:
      type: photo
      prompt: Take photo of clear water and chemical test strip
\end{minted}

% ─────────────────────────────────────────────────────────────────────────────
\section{The Delegation Chain}
\label{sec:concepts-chain}
% ─────────────────────────────────────────────────────────────────────────────

Away-mode delegation follows a simple activation and expiry chain.

\paragraph{How it works.}
\begin{enumerate}
  \item At each routine execution, the runtime calls \texttt{IsAwayModeActive()} with the
        current timestamp.
  \item If a \texttt{delegations} entry is active (current date is between \texttt{active.from}
        and \texttt{active.to}), the runtime checks whether the routine or its zone appears in
        the \texttt{assigns} list.
  \item If matched, the delegate's identity is substituted as executor for that routine run.
  \item If the routine is not in any active delegation's scope, it falls back to the homeowner.
  \item When all active delegations expire, the homeowner resumes as the default executor
        automatically with no configuration change required.
\end{enumerate}

\paragraph{Overlapping delegations.}
Multiple \texttt{delegations} entries may be active simultaneously. If more than one active
delegation covers the same routine or zone, the first matching entry in the list wins.

% ─────────────────────────────────────────────────────────────────────────────
\section{SLA / Cadence Concepts}
\label{sec:concepts-sla}
% ─────────────────────────────────────────────────────────────────────────────

\texttt{gert.home} supports three cadence models for scheduling recurring routines and
replacement of consumables.

\subsection{Fixed Interval}

The most common cadence. The routine is due every \textit{N} days or weeks after its last
completion.

\begin{minted}{yaml}
cadence:
  every: 7d   # due every 7 days
\end{minted}

Supported units: \texttt{d} (days), \texttt{w} (weeks).

\subsection{Seasonal Schedule}

For tasks whose frequency varies by season (e.g., lawn care). Each season key maps to a
fixed interval that applies while that season is current.

\begin{minted}{yaml}
cadence:
  seasonal:
    summer: 7d
    autumn: 10d
    winter: 21d
    spring: 10d
\end{minted}

Season boundaries are determined by the hemisphere setting derived from
\texttt{GERT\_HOME\_TZ} or defaulting to UTC (northern hemisphere seasons).

\subsection{Consumable Replacement}

Consumables use a \texttt{replace\_every} field rather than a routine cadence.

\begin{minted}{yaml}
consumables:
  - id: pool_filter_sand
    replace_every: 180d
    last_replaced: 2024-11-01
\end{minted}

\subsection{Notification Lead Times}

Each routine may declare reminder and escalation windows:

\begin{minted}{yaml}
notifications:
  remind_hours_before: 4
  escalate_hours_overdue: 24
\end{minted}

\begin{description}
  \item[\texttt{remind\_hours\_before}] How many hours before the due date to send a reminder
        to the executor.
  \item[\texttt{escalate\_hours\_overdue}] How many hours past the due date before an overdue
        alert is sent to the homeowner.
\end{description}

% ─────────────────────────────────────────────────────────────────────────────
\section{Delegation Transfer}
\label{sec:concepts-transfer}
% ─────────────────────────────────────────────────────────────────────────────

Delegation in \texttt{gert.home} is always explicit: a \texttt{delegations} block with a
defined \texttt{from}/\texttt{to} date range. There is no implicit or ad-hoc transfer.

\paragraph{Explicit delegation.}
The homeowner authors a \texttt{delegations} entry before leaving. The entry specifies the
delegate's identity, the active window, the scope of assigned tasks, and the permissions
granted.

\paragraph{Automatic expiry.}
When the current date passes \texttt{active.to}, the delegation expires without any action
required. The runtime's \texttt{IsAwayModeActive()} call returns \texttt{false} and the
homeowner resumes as executor.

\paragraph{Permissions scope.}
Each delegation carries an explicit permissions map. This prevents accidental scope creep:
a delegate with \texttt{can\_modify\_routines: false} cannot alter the cadence or evidence
requirements of any routine, even those they are assigned to execute.

\chapter{Schema Reference}
\label{ch:schema-reference}

This chapter is the authoritative reference for every schema construct in the
\texttt{gert.home} Domain Kit. It covers the Kit declaration, the property metadata block,
all step types, evidence types, and consumables. For the core gert schema (step types
\texttt{cli}, \texttt{manual}, \texttt{tool}, \texttt{branch}, \texttt{iterate}), refer to
the \emph{gert Design Document}, §3.

% ─────────────────────────────────────────────────────────────────────────────
\section{Kit Declaration and \texttt{apiVersion}}
\label{sec:schema-apiversion}
% ─────────────────────────────────────────────────────────────────────────────

A \texttt{gert.home} property file declares both the core schema version and the Kit at the
top of the file. The file extension is \texttt{.home.yaml}.

\begin{minted}{yaml}
apiVersion: runbook/v1
kit: home/v0

property:
  name: Casa Santiago
  zones:
    - id: pool
      label: Swimming Pool
      description: Saltwater pool with heat pump and LED lighting
\end{minted}

\begin{description}
  \item[\texttt{apiVersion: runbook/v1}] Required. Selects the gert core schema.
  \item[\texttt{kit: home/v0}] Required. Activates the \texttt{gert.home} Kit compiler,
        making \texttt{home.*} step types and property metadata fields available.
\end{description}

% ─────────────────────────────────────────────────────────────────────────────
\section{Property Metadata Fields}
\label{sec:schema-metadata}
% ─────────────────────────────────────────────────────────────────────────────

\subsection{\texttt{property} Block}

The \texttt{property} block declares the household unit identity and its named zones.

\begin{minted}{yaml}
property:
  name: Casa Santiago
  zones:
    - id: pool
      label: Swimming Pool
      description: Saltwater pool with heat pump and LED lighting
    - id: front_lawn
      label: Front Lawn
\end{minted}

\begin{description}
  \item[\texttt{property.name}] Display name for the property. Shown in reports and
        notifications.
  \item[\texttt{property.zones}] List of named areas within the property (see below).
\end{description}

\subsection{\texttt{zones[]} Fields}

\begin{description}
  \item[\texttt{id}] Unique identifier for the zone. Referenced by routines,
        incident templates, and delegations.
  \item[\texttt{label}] Human-readable display name.
  \item[\texttt{description}] Optional. Free-text description of the zone.
  \item[\texttt{metadata}] Optional. Arbitrary key-value pairs for additional context.
\end{description}

\subsection{\texttt{assets[]} Fields}

Assets are tracked physical items within a zone.

\begin{minted}{yaml}
assets:
  - id: pool_pump
    zone: pool
    label: Hayward SuperPump VS
    installed: 2023-06-15
    warranty_expires: 2025-06-15
    metadata:
      model: VS-100
\end{minted}

\begin{description}
  \item[\texttt{id}] Unique identifier for the asset. Referenced by routines and consumables.
  \item[\texttt{zone}] Zone ID where the asset is located.
  \item[\texttt{label}] Human-readable display name.
  \item[\texttt{installed}] ISO 8601 date of installation (\texttt{YYYY-MM-DD}).
  \item[\texttt{warranty\_expires}] Optional. ISO 8601 date the warranty expires.
  \item[\texttt{metadata}] Optional. Arbitrary key-value pairs (e.g., model number, serial).
\end{description}

% ─────────────────────────────────────────────────────────────────────────────
\section{Step Types}
\label{sec:schema-step-types}
% ─────────────────────────────────────────────────────────────────────────────

The \texttt{gert.home} kit defines three step types. All are lowered to core primitives at
compile time.

\subsection{\texttt{home.routine}}

A scheduled recurring maintenance task. Lowers to an \texttt{iterate} core step containing
a \texttt{manual} step, one execution per due routine instance.

\begin{minted}{yaml}
routines:
  - id: pool_clean
    label: Pool cleaning
    description: Weekly pool water quality check and surface clean
    zone: pool
    cadence:
      every: 7d
    evidence:
      type: photo
      prompt: Take photo of clear water and chemical test strip
    executor:
      role: gardener
    notifications:
      remind_hours_before: 4
      escalate_hours_overdue: 24
\end{minted}

\begin{description}
  \item[\texttt{id}] Required. Unique identifier for the routine.
  \item[\texttt{label}] Required. Short display name.
  \item[\texttt{description}] Optional. Longer description of the task.
  \item[\texttt{zone}] Conditional. Zone ID scoping this routine. One of \texttt{zone} or
        \texttt{asset} must be provided.
  \item[\texttt{asset}] Conditional. Asset ID scoping this routine. One of \texttt{zone} or
        \texttt{asset} must be provided.
  \item[\texttt{cadence.every}] Conditional. Fixed interval (e.g., \texttt{7d}, \texttt{2w}).
  \item[\texttt{cadence.seasonal}] Conditional. Season-keyed intervals
        (\texttt{summer}, \texttt{autumn}, \texttt{winter}, \texttt{spring}).
  \item[\texttt{evidence.type}] Optional. Evidence type required on completion.
  \item[\texttt{evidence.prompt}] Optional. Instruction shown to the executor.
  \item[\texttt{executor.role}] Optional. Advisory role hint (e.g., \texttt{gardener},
        \texttt{plumber}). Not enforced at runtime in \texttt{home/v0}.
  \item[\texttt{notifications.remind\_hours\_before}] Optional. Hours before due date to
        send reminder.
  \item[\texttt{notifications.escalate\_hours\_overdue}] Optional. Hours past due date to
        send overdue alert.
\end{description}

\subsection{\texttt{home.incident}}

A reactive repair workflow with sequential, parallel, and branching steps. Each incident
step type lowers to a core primitive: \texttt{human\_task} → \texttt{manual},
\texttt{decision} → \texttt{branch}, \texttt{parallel} → \texttt{parallel}.

\begin{minted}{yaml}
incident_templates:
  - id: leak
    label: Water leak
    description: Emergency water leak response procedure
    zones: [pool, backyard, garage]
    steps:
      - id: assess
        type: human_task
        label: Locate and assess the leak
        evidence:
          type: photo
          prompt: Take photo of leak source and affected area
      - id: shutoff
        type: human_task
        label: Shut off water supply
        depends_on: [assess]
        evidence:
          type: photo
      - id: fix
        type: human_task
        label: Perform repair
        depends_on: [shutoff]
        evidence:
          type: photo
\end{minted}

\begin{description}
  \item[\texttt{id}] Required. Unique identifier for the incident template.
  \item[\texttt{label}] Required. Short display name.
  \item[\texttt{description}] Optional. Longer description.
  \item[\texttt{zones}] Optional. List of zone IDs this incident may affect.
  \item[\texttt{assets}] Optional. List of asset IDs this incident may affect.
  \item[\texttt{steps}] Required. List of IncidentStep objects (see below).
\end{description}

\paragraph{IncidentStep types.}

\begin{description}
  \item[\texttt{human\_task}] A manual procedure requiring human action. Fields:
        \texttt{id}, \texttt{type}, \texttt{label}, \texttt{description}, \texttt{evidence},
        \texttt{depends\_on}.
  \item[\texttt{decision}] A branching choice. Fields: \texttt{id}, \texttt{type},
        \texttt{label}, \texttt{choices} (list of \texttt{\{value, label, next\_step\}}).
  \item[\texttt{parallel}] Sub-steps executed concurrently. Fields: \texttt{id},
        \texttt{type}, \texttt{label}, \texttt{parallel\_steps}.
\end{description}

\subsection{\texttt{home.delegate}}

An away-mode delegation block. Checked at run-time via \texttt{IsAwayModeActive()};
when active, substitutes the delegate identity as executor for assigned routines or zones.

\begin{minted}{yaml}
delegations:
  - delegate:
      name: Carlos Méndez
      contact: "+56 9 8765 4321"
      email: carlos@example.com
    active:
      from: 2025-04-25
      to: 2025-05-02
    assigns:
      - routine: water_plants
      - zone: pool
    permissions:
      can_report_incidents: true
      can_modify_routines: false
      can_view_history: true
    notifications:
      remind_delegate_hours_before: 24
      notify_owner_on_completion: true
      notify_owner_if_overdue: true
\end{minted}

\begin{description}
  \item[\texttt{delegate.name}] Required. Display name of the delegate.
  \item[\texttt{delegate.contact}] Required. Phone number or contact method.
  \item[\texttt{delegate.email}] Optional. Email address.
  \item[\texttt{active.from}] Required. Delegation start date (\texttt{YYYY-MM-DD}).
  \item[\texttt{active.to}] Required. Delegation end date (\texttt{YYYY-MM-DD}).
  \item[\texttt{assigns}] Required. List of \texttt{\{routine: id\}} or
        \texttt{\{zone: id\}} entries scoping the delegation.
  \item[\texttt{permissions.can\_report\_incidents}] Optional. Default \texttt{false}.
  \item[\texttt{permissions.can\_modify\_routines}] Optional. Default \texttt{false}.
  \item[\texttt{permissions.can\_view\_history}] Optional. Default \texttt{false}.
\end{description}

% ─────────────────────────────────────────────────────────────────────────────
\section{Evidence Types}
\label{sec:schema-evidence}
% ─────────────────────────────────────────────────────────────────────────────

All evidence types share two common fields:

\begin{description}
  \item[\texttt{type}] Required. One of: \texttt{photo}, \texttt{note},
        \texttt{checklist}, \texttt{none}.
  \item[\texttt{prompt}] Optional. Instruction or question shown to the executor at the
        time of evidence submission.
\end{description}

\subsection{\texttt{photo}}

A photograph proving task completion. Required for most outdoor and physical tasks.

\begin{minted}{yaml}
evidence:
  type: photo
  prompt: Take photo of clear water and chemical test strip
\end{minted}

\subsection{\texttt{note}}

A free-text note from the executor. Useful for recording cost, measurements, or observations.

\begin{minted}{yaml}
evidence:
  type: note
  prompt: Record the chemical readings and any materials used
\end{minted}

\subsection{\texttt{checklist}}

A structured checklist that must be completed item by item before the step closes.

\begin{minted}{yaml}
evidence:
  type: checklist
  prompt: Confirm all safety checks before closing the incident
\end{minted}

\subsection{\texttt{none}}

No evidence required. The step is self-certifying on completion.

\begin{minted}{yaml}
evidence:
  type: none
\end{minted}

% ─────────────────────────────────────────────────────────────────────────────
\section{Consumables}
\label{sec:schema-consumables}
% ─────────────────────────────────────────────────────────────────────────────

Consumables are items on a replacement schedule with optional stock tracking.

\begin{minted}{yaml}
consumables:
  - id: pool_filter_sand
    label: Pool sand filter media
    zone: pool
    asset: pool_pump
    replace_every: 180d
    last_replaced: 2024-11-01
    stock:
      current: 1
      reorder_at: 1
    triggers_routine: pool_filter_replace
\end{minted}

\begin{description}
  \item[\texttt{id}] Required. Unique identifier.
  \item[\texttt{label}] Required. Display name.
  \item[\texttt{description}] Optional. Longer description.
  \item[\texttt{zone}] Optional. Zone where the consumable is used.
  \item[\texttt{asset}] Optional. Asset ID this consumable belongs to.
  \item[\texttt{replace\_every}] Required. Replacement interval (e.g., \texttt{180d}).
  \item[\texttt{last\_replaced}] Required. Date of last replacement (\texttt{YYYY-MM-DD}).
  \item[\texttt{stock.current}] Optional. Current stock quantity on hand.
  \item[\texttt{stock.reorder\_at}] Optional. Stock level that triggers a reorder alert.
  \item[\texttt{triggers\_routine}] Optional. Routine ID to fire when replacement is due.
\end{description}

\chapter{Authoring Guide}
\label{ch:authoring-guide}

This chapter teaches you to write a \texttt{.home.yaml} property file from scratch. It walks
through zone and asset declaration, routine authoring, incident templates, delegation
configuration, and consumable tracking, ending with a complete pool-maintenance example.

% ─────────────────────────────────────────────────────────────────────────────
\section{Starting a New Property File}
\label{sec:authoring-start}
% ─────────────────────────────────────────────────────────────────────────────

Every \texttt{gert.home} property file begins with the Kit declaration and a minimal
\texttt{property} block.

\begin{minted}{yaml}
apiVersion: runbook/v1
kit: home/v0

property:
  name: My Home
  zones: []

assets: []
routines: []
incident_templates: []
consumables: []
delegations: []
\end{minted}

Start by naming the property and identifying the major areas (zones) that require maintenance.
Assets, routines, and incidents can be added incrementally. A property file with only zones
and assets is valid and will compile without error.

% ─────────────────────────────────────────────────────────────────────────────
\section{Declaring Zones and Assets}
\label{sec:authoring-zones-assets}
% ─────────────────────────────────────────────────────────────────────────────

\paragraph{Zones} are named areas of the property. Use zones for locations that contain
multiple assets or are referenced by multiple routines (e.g., \texttt{pool},
\texttt{front\_lawn}, \texttt{garage}).

\paragraph{Assets} are specific tracked physical items. Use assets when you need to record
installation date, warranty expiry, or model metadata for a particular device or appliance.

\begin{minted}{yaml}
property:
  name: Casa Santiago
  zones:
    - id: pool
      label: Swimming Pool
      description: Saltwater pool with heat pump and LED lighting
    - id: front_lawn
      label: Front Lawn

assets:
  - id: pool_pump
    zone: pool
    label: Hayward SuperPump VS
    installed: 2023-06-15
    warranty_expires: 2025-06-15
    metadata:
      model: VS-100
\end{minted}

\begin{tcolorbox}[colback=yellow!10,colframe=orange!70,title=Rule of Thumb]
  Use a \textbf{zone} for any area you refer to by location (``the pool area'', ``the
  garage''). Use an \textbf{asset} for any specific item you might need to look up by model
  number, warranty, or service history.
\end{tcolorbox}

% ─────────────────────────────────────────────────────────────────────────────
\section{Defining Routines}
\label{sec:authoring-routines}
% ─────────────────────────────────────────────────────────────────────────────

Routines are the core of the property file. Each routine declares a recurring task with a
cadence, an evidence requirement, and optional notifications.

\begin{minted}{yaml}
routines:
  - id: pool_clean
    label: Pool cleaning
    zone: pool
    cadence:
      every: 7d
    evidence:
      type: photo
      prompt: Take photo of clear water and chemical test strip
    notifications:
      remind_hours_before: 4
      escalate_hours_overdue: 24

  - id: lawn_mow
    label: Mow front lawn
    zone: front_lawn
    cadence:
      seasonal:
        summer: 7d
        autumn: 10d
        winter: 21d
        spring: 10d
    evidence:
      type: photo
    executor:
      role: gardener
\end{minted}

\paragraph{Fixed vs.\ seasonal cadence.}
Use \texttt{every: Nd} when the task frequency does not change with the seasons (e.g., pool
cleaning, bin collection). Use \texttt{seasonal} when frequency varies significantly across
the year (e.g., lawn mowing, irrigation). Mixing both cadence types in the same routine
is not allowed --- choose one.

% ─────────────────────────────────────────────────────────────────────────────
\section{Using \texttt{home.incident}}
\label{sec:authoring-incidents}
% ─────────────────────────────────────────────────────────────────────────────

Use \texttt{incident\_templates} for reactive, event-driven workflows --- situations you
cannot schedule in advance. Common examples: water leaks, appliance failures, power outages.

\begin{minted}{yaml}
incident_templates:
  - id: leak
    label: Water leak
    zones: [pool, backyard, garage]
    steps:
      - id: assess
        type: human_task
        label: Locate and assess the leak
        evidence:
          type: photo
          prompt: Take photo of leak source and affected area
      - id: shutoff
        type: human_task
        label: Shut off water supply
        depends_on: [assess]
        evidence:
          type: photo
      - id: severity_check
        type: decision
        label: Can you fix it yourself?
        depends_on: [shutoff]
        choices:
          - value: yes
            label: Yes, I can fix it
            next_step: fix
          - value: no
            label: No, call a plumber
            next_step: call_plumber
      - id: fix
        type: human_task
        label: Perform repair
        evidence:
          type: photo
      - id: call_plumber
        type: human_task
        label: Contact a licensed plumber
        evidence:
          type: note
          prompt: Record plumber name, contact number, and estimated arrival time
\end{minted}

\paragraph{Structuring steps.}
Use \texttt{depends\_on} to chain steps sequentially. Steps without \texttt{depends\_on}
may start immediately when the incident is opened. Use \texttt{type: decision} for branching
choices and \texttt{type: parallel} for tasks that can be done concurrently.

% ─────────────────────────────────────────────────────────────────────────────
\section{Configuring Delegation}
\label{sec:authoring-delegation}
% ─────────────────────────────────────────────────────────────────────────────

Add a \texttt{delegations} entry before you leave the property. Specify the delegate's
contact details, the exact date range, which routines or zones they are assigned to, and
what permissions they have.

\begin{minted}{yaml}
delegations:
  - delegate:
      name: Carlos Méndez
      contact: "+56 9 8765 4321"
      email: carlos@example.com
    active:
      from: 2025-04-25
      to: 2025-05-02
    assigns:
      - routine: water_plants
      - zone: pool
    permissions:
      can_report_incidents: true
      can_modify_routines: false
      can_view_history: true
    notifications:
      remind_delegate_hours_before: 24
      notify_owner_on_completion: true
\end{minted}

\paragraph{Timing the away window.}
Set \texttt{active.from} to the day you leave and \texttt{active.to} to the day you return.
The delegation activates at midnight on \texttt{active.from} and expires at midnight on
\texttt{active.to} (in the timezone set by \texttt{GERT\_HOME\_TZ}).

\paragraph{Scoping assigns.}
Use \texttt{\{routine: id\}} to assign specific routines, and \texttt{\{zone: id\}} to
assign all routines whose \texttt{zone} field matches. Zone assignment is useful when you want
the delegate to handle everything in a given area without listing each routine individually.

% ─────────────────────────────────────────────────────────────────────────────
\section{Tracking Consumables}
\label{sec:authoring-consumables}
% ─────────────────────────────────────────────────────────────────────────────

Consumables track items that must be replaced on a schedule. The runtime computes the next
replacement date from \texttt{last\_replaced} + \texttt{replace\_every} and alerts when due.

\begin{minted}{yaml}
consumables:
  - id: pool_filter_sand
    label: Pool sand filter media
    zone: pool
    asset: pool_pump
    replace_every: 180d
    last_replaced: 2024-11-01
    stock:
      current: 1
      reorder_at: 1
    triggers_routine: pool_filter_replace
\end{minted}

After replacing a consumable, update \texttt{last\_replaced} to today's date. If
\texttt{triggers\_routine} is set, the runtime will fire the named routine when the
consumable's replacement date arrives.

% ─────────────────────────────────────────────────────────────────────────────
\section{Evidence Best Practices}
\label{sec:authoring-evidence}
% ─────────────────────────────────────────────────────────────────────────────

\begin{itemize}
  \item Use \texttt{photo} for outdoor and physical tasks where completion is visible
        (mowing, pool cleaning, leak repair). Photos provide tamper-evident proof and are
        SHA-256 hashed in the run record.
  \item Use \texttt{note} for tasks involving cost, measurements, or contractor
        communications (e.g., recording chemical readings, logging repair costs).
  \item Use \texttt{checklist} for multi-step verification tasks where partial completion
        must be tracked (e.g., pre-departure safety checks, appliance service checklists).
  \item Use \texttt{none} sparingly --- only for tasks where completion is self-evident and
        no record is required (e.g., resetting a timer, clearing a filter alarm).
  \item Always write a \texttt{prompt} for \texttt{photo} and \texttt{note} evidence. A
        clear prompt reduces ambiguous submissions and improves the audit trail.
  \item For incident steps, require \texttt{photo} evidence on the final \texttt{human\_task}
        step so you have documented proof of the completed repair.
\end{itemize}

% ─────────────────────────────────────────────────────────────────────────────
\section{Complete Example: Pool Maintenance Property}
\label{sec:authoring-complete-example}
% ─────────────────────────────────────────────────────────────────────────────

The following is a complete \texttt{.home.yaml} file for a property with a swimming pool and
front lawn. It demonstrates all major features: zones, assets, fixed and seasonal routines,
an incident template, a consumable, and a delegation.

\begin{minted}{yaml}
apiVersion: runbook/v1
kit: home/v0

property:
  name: Casa Santiago
  zones:
    - id: pool
      label: Swimming Pool
      description: Saltwater pool with heat pump and LED lighting
    - id: front_lawn
      label: Front Lawn

assets:
  - id: pool_pump
    zone: pool
    label: Hayward SuperPump VS
    installed: 2023-06-15
    warranty_expires: 2025-06-15
    metadata:
      model: VS-100

routines:
  - id: pool_clean
    label: Pool cleaning
    zone: pool
    cadence:
      every: 7d
    evidence:
      type: photo
      prompt: Take photo of clear water and chemical test strip
    notifications:
      remind_hours_before: 4
      escalate_hours_overdue: 24

  - id: lawn_mow
    label: Mow front lawn
    zone: front_lawn
    cadence:
      seasonal:
        summer: 7d
        autumn: 10d
        winter: 21d
        spring: 10d
    evidence:
      type: photo
    executor:
      role: gardener

incident_templates:
  - id: leak
    label: Water leak
    zones: [pool, backyard, garage]
    steps:
      - id: assess
        type: human_task
        label: Locate and assess the leak
        evidence:
          type: photo
          prompt: Take photo of leak source and affected area
      - id: shutoff
        type: human_task
        label: Shut off water supply
        depends_on: [assess]
        evidence:
          type: photo
      - id: fix
        type: human_task
        label: Perform repair
        depends_on: [shutoff]
        evidence:
          type: photo

consumables:
  - id: pool_filter_sand
    label: Pool sand filter media
    zone: pool
    asset: pool_pump
    replace_every: 180d
    last_replaced: 2024-11-01
    stock:
      current: 1
      reorder_at: 1

delegations:
  - delegate:
      name: Carlos Méndez
      contact: "+56 9 8765 4321"
    active:
      from: 2025-04-25
      to: 2025-05-02
    assigns:
      - routine: pool_clean
      - zone: front_lawn
    permissions:
      can_report_incidents: true
      can_modify_routines: false
      can_view_history: true
    notifications:
      remind_delegate_hours_before: 24
      notify_owner_on_completion: true
\end{minted}

\chapter{gert Integration}
\label{ch:gert-integration}

This chapter describes how \texttt{gert.home} integrates with the gert runtime: how
property files are parsed, validated, and lowered to core primitives; how away mode is
enforced; how cadence checking works; and what trace events are emitted during execution.

% ─────────────────────────────────────────────────────────────────────────────
\section{Core Integration Hooks}
\label{sec:integration-hooks}
% ─────────────────────────────────────────────────────────────────────────────

A \texttt{.home.yaml} property file passes through four stages before any step executes:

\begin{enumerate}
  \item \textbf{Parse} --- The YAML file is parsed and the \texttt{apiVersion} and
        \texttt{kit} fields are read. The \texttt{home/v0} Kit compiler is activated.
  \item \textbf{Validate} --- The property file is validated against the \texttt{gert.home}
        JSON Schema. Zone and asset references are resolved; missing zone IDs or invalid
        cadence formats produce compile errors.
  \item \textbf{Lower} --- \texttt{home.*} step types are lowered to core gert primitives:
        \texttt{home.routine} → \texttt{iterate} + \texttt{manual}; incident
        \texttt{human\_task} → \texttt{manual}; \texttt{decision} → \texttt{branch};
        \texttt{parallel} → \texttt{parallel}.
  \item \textbf{Execute} --- The gert runtime executes the lowered runbook, checking
        cadence due dates, calling \texttt{IsAwayModeActive()}, and emitting trace events at
        each step boundary.
\end{enumerate}

The gert runtime never sees \texttt{home.*} vocabulary after the lower stage. This means
all core gert tooling (trace readers, evidence bundlers, CI integrations) works with
\texttt{gert.home} runs without modification.

% ─────────────────────────────────────────────────────────────────────────────
\section{Role-Based Gating}
\label{sec:integration-role-gating}
% ─────────────────────────────────────────────────────────────────────────────

\subsection{Executor Role Hints}

The \texttt{executor.role} field on a \texttt{home.routine} is advisory in
\texttt{home/v0}. The runtime reads the field and includes it in notification messages and
the task UI, but does not block execution if the current actor does not match the hint.

\begin{tcolorbox}[colback=yellow!10,colframe=orange!70,title=Note]
  \texttt{home/v0} does not enforce executor role hints at runtime. A homeowner or active
  delegate can always execute any step regardless of the \texttt{executor.role} value. Hard
  role gates are planned for \texttt{home/v1}.
\end{tcolorbox}

\subsection{Delegation Scope Enforcement}

The \texttt{assigns} block in a \texttt{delegations} entry is enforced at runtime.
When \texttt{IsAwayModeActive()} returns \texttt{true}:

\begin{itemize}
  \item The runtime checks whether the current routine's ID appears in any active
        delegation's \texttt{assigns} list as \texttt{\{routine: id\}}.
  \item The runtime also checks whether the routine's \texttt{zone} ID appears as
        \texttt{\{zone: id\}}.
  \item If neither match, the routine is not delegated and the homeowner remains executor.
  \item If a match is found, the delegate's name and contact are substituted as executor
        identity for that run.
\end{itemize}

% ─────────────────────────────────────────────────────────────────────────────
\section{Command / Action Scope}
\label{sec:integration-commands}
% ─────────────────────────────────────────────────────────────────────────────

\texttt{gert.home} routines and incidents generate \texttt{manual} steps --- human-action
steps that pause execution until a person confirms completion and submits evidence.
The kit does not generate \texttt{cli} steps. There are no shell commands, no environment
variable access, and no network calls in a \texttt{gert.home} execution. This keeps the
security surface minimal for household contexts.

% ─────────────────────────────────────────────────────────────────────────────
\section{Away Mode Enforcement}
\label{sec:integration-away-mode}
% ─────────────────────────────────────────────────────────────────────────────

Away mode is evaluated at the start of each routine execution via the
\texttt{IsAwayModeActive()} runtime function:

\begin{enumerate}
  \item The runtime reads all \texttt{delegations} entries from the property file.
  \item It compares the current timestamp against each entry's \texttt{active.from} and
        \texttt{active.to} fields using the timezone from \texttt{GERT\_HOME\_TZ}.
  \item If one or more entries are active, \texttt{IsAwayModeActive()} returns \texttt{true}
        and the runtime proceeds with delegation scope checks.
  \item If no entries are active, \texttt{IsAwayModeActive()} returns \texttt{false} and
        the homeowner is executor for all routines.
\end{enumerate}

When away mode is active, a \texttt{delegation\_active} trace event is emitted for each
routine execution that matches an active delegation scope.

% ─────────────────────────────────────────────────────────────────────────────
\section{Cadence Enforcement}
\label{sec:integration-cadence}
% ─────────────────────────────────────────────────────────────────────────────

Before firing a \texttt{home.routine} step, the runtime checks whether the routine is due:

\begin{enumerate}
  \item The runtime reads the last completion timestamp from the run directory
        (\texttt{runs/\{routine-id\}/latest/trace.jsonl}).
  \item It computes the next due date: last completion + cadence interval (or the seasonal
        interval for the current date).
  \item If the current timestamp is before the due date, the routine is skipped and a
        \texttt{routine\_skipped} trace event is emitted.
  \item If the current timestamp is within \texttt{notifications.remind\_hours\_before} of
        the due date, a reminder notification is sent to the executor.
  \item If the current timestamp is past \texttt{due date + notifications.escalate\_hours\_overdue},
        an overdue alert is sent to the homeowner regardless of delegation state.
\end{enumerate}

% ─────────────────────────────────────────────────────────────────────────────
\section{The Audit Trail}
\label{sec:integration-audit-trail}
% ─────────────────────────────────────────────────────────────────────────────

The \texttt{gert.home} runtime emits structured trace events throughout execution. These are
written to \texttt{runs/\{routine-id\}/\{date\}/trace.jsonl} in newline-delimited JSON.

\begin{description}
  \item[\texttt{routine\_due}] Emitted when a routine's due date arrives. Fields:
        \texttt{routine\_id}, \texttt{zone}, \texttt{executor}.
  \item[\texttt{routine\_completed}] Emitted when all steps of a routine run are completed.
        Fields: \texttt{routine\_id}, \texttt{duration\_minutes}.
  \item[\texttt{incident\_opened}] Emitted when an incident template is instantiated.
        Fields: \texttt{incident\_id}, \texttt{opened\_by}.
  \item[\texttt{incident\_step\_completed}] Emitted when each incident step finishes.
        Fields: \texttt{incident\_id}, \texttt{step\_id}, \texttt{type}.
  \item[\texttt{evidence\_submitted}] Emitted when evidence is captured. Fields:
        \texttt{step\_id}, \texttt{evidence\_type}, \texttt{submitted\_by}, \texttt{sha256}.
  \item[\texttt{delegation\_active}] Emitted at the start of any routine execution where
        an active delegation scope matches. Fields: \texttt{delegate\_name},
        \texttt{routine\_id}, \texttt{active\_until}.
\end{description}

% ─────────────────────────────────────────────────────────────────────────────
\section{Complete Integration Example}
\label{sec:integration-complete-example}
% ─────────────────────────────────────────────────────────────────────────────

The following shows a property file with an active delegation alongside the expected trace
output for a pool-cleaning routine run during the away window.

\begin{minted}{yaml}
apiVersion: runbook/v1
kit: home/v0

property:
  name: Casa Santiago
  zones:
    - id: pool
      label: Swimming Pool

routines:
  - id: pool_clean
    label: Pool cleaning
    zone: pool
    cadence:
      every: 7d
    evidence:
      type: photo
      prompt: Take photo of clear water and chemical test strip
    notifications:
      remind_hours_before: 4
      escalate_hours_overdue: 24

delegations:
  - delegate:
      name: Carlos Méndez
      contact: "+56 9 8765 4321"
    active:
      from: 2025-04-25
      to: 2025-05-02
    assigns:
      - zone: pool
    permissions:
      can_report_incidents: true
      can_modify_routines: false
      can_view_history: true
\end{minted}

Trace events emitted on 2025-04-28 when the routine fires:

\begin{minted}{yaml}
{"seq":1,"ts":"2025-04-28T08:00:00Z","type":"delegation_active",
 "delegate_name":"Carlos Méndez","routine_id":"pool_clean",
 "active_until":"2025-05-02"}
{"seq":2,"ts":"2025-04-28T08:00:01Z","type":"routine_due",
 "routine_id":"pool_clean","zone":"pool","executor":"Carlos Méndez"}
{"seq":3,"ts":"2025-04-28T09:15:00Z","type":"evidence_submitted",
 "step_id":"pool_clean","evidence_type":"photo",
 "submitted_by":"Carlos Méndez",
 "sha256":"a3f2...c1d9"}
{"seq":4,"ts":"2025-04-28T09:15:01Z","type":"routine_completed",
 "routine_id":"pool_clean","duration_minutes":75}
\end{minted}

\chapter{Evidence and Tracing}
\label{ch:evidence-tracing}

This chapter describes the evidence capture model used by \texttt{gert.home}: what is
captured, when, how it is stored, and how to verify and read the evidence bundle from a
completed run. The underlying trace format is defined by gert core (\emph{gert Design
Document}, §12); this chapter focuses on the \texttt{gert.home}-layer semantics.

% ─────────────────────────────────────────────────────────────────────────────
\section{The Evidence Model}
\label{sec:evidence-model}
% ─────────────────────────────────────────────────────────────────────────────

\texttt{gert.home} captures evidence at the step level. Each piece of evidence is associated
with the step that required it, the executor who submitted it, and the timestamp of submission.

\paragraph{What gets captured.}
\begin{itemize}
  \item \textbf{Photos} --- Binary image file stored under the run directory.
        Captured: filename, timestamp, zone and asset context, SHA-256 hash.
  \item \textbf{Notes} --- Free-text string submitted by the executor.
        Captured: note text, executor identity (name or delegate name), timestamp.
  \item \textbf{Checklists} --- Completion status for each item in the checklist.
        Captured: item label, checked/unchecked, timestamp per item, overall completion.
\end{itemize}

\paragraph{What is NOT evidence.}
The following events are recorded in the trace log but are \emph{not} included in the
evidence bundle:
\begin{itemize}
  \item Reminder notifications
  \item Overdue escalation alerts
  \item Scheduling events (\texttt{routine\_due}, \texttt{routine\_skipped})
  \item Delegation activation events
\end{itemize}

% ─────────────────────────────────────────────────────────────────────────────
\section{Evidence Record Structure}
\label{sec:evidence-record}
% ─────────────────────────────────────────────────────────────────────────────

Each evidence submission is recorded as a structured object in the trace log.

\begin{minted}{yaml}
type: evidence_submitted
seq: 3
ts: "2025-04-28T09:15:00Z"
step_id: pool_clean
evidence_type: photo
label: Pool water and test strip
submitted_by: Carlos Méndez
role: delegate
content:
  filename: pool_clean_2025-04-28.jpg
  sha256: a3f2b19e4c7d8e6f1a2b3c4d5e6f7a8b9c0d1e2f3a4b5c6d7e8f9a0b1c2d3e4f5
\end{minted}

\begin{description}
  \item[\texttt{type}] Always \texttt{evidence\_submitted} for evidence records.
  \item[\texttt{seq}] Sequence number within the run's trace log.
  \item[\texttt{ts}] ISO 8601 timestamp of submission.
  \item[\texttt{step\_id}] ID of the step that required this evidence.
  \item[\texttt{evidence\_type}] One of \texttt{photo}, \texttt{note}, \texttt{checklist}.
  \item[\texttt{label}] Human-readable label from the routine or incident step definition.
  \item[\texttt{submitted\_by}] Name of the person who submitted the evidence.
  \item[\texttt{role}] Either \texttt{homeowner} or \texttt{delegate}.
  \item[\texttt{content}] Type-specific payload. For photos: filename and SHA-256 hash.
        For notes: free text. For checklists: list of item completion records.
\end{description}

% ─────────────────────────────────────────────────────────────────────────────
\section{Run Directory Structure}
\label{sec:evidence-run-dir}
% ─────────────────────────────────────────────────────────────────────────────

Each routine run is stored in a directory under the configured run root
(\texttt{GERT\_HOME\_RUN\_DIR}, default: \texttt{./runs}).

\begin{minted}{yaml}
runs/
  pool_clean/
    2025-04-28/
      trace.jsonl          # Newline-delimited JSON trace events
      evidence/
        pool_clean_2025-04-28.jpg     # Photo evidence files
        pool_clean_2025-04-28.sha256  # SHA-256 sidecar file
    2025-05-05/
      trace.jsonl
      evidence/
        pool_clean_2025-05-05.jpg
        pool_clean_2025-05-05.sha256
    latest -> 2025-05-05   # Symlink to most recent run
  leak/
    2025-04-29/
      trace.jsonl
      evidence/
        assess_2025-04-29.jpg
        shutoff_2025-04-29.jpg
        fix_2025-04-29.jpg
\end{minted}

Each run date directory is named \texttt{YYYY-MM-DD} using the timezone from
\texttt{GERT\_HOME\_TZ}. The \texttt{latest} symlink always points to the most recent run
for quick access.

% ─────────────────────────────────────────────────────────────────────────────
\section{Reading the Timeline Projection}
\label{sec:evidence-timeline}
% ─────────────────────────────────────────────────────────────────────────────

The \texttt{trace.jsonl} file contains one JSON object per line, ordered by sequence number.
Key fields in each event:

\begin{description}
  \item[\texttt{seq}] Monotonically increasing sequence number within the run. Useful for
        ordering events when timestamps have the same second.
  \item[\texttt{ts}] ISO 8601 timestamp (UTC).
  \item[\texttt{step\_id}] ID of the step that emitted this event.
  \item[\texttt{type}] Event type. Home-layer types: \texttt{routine\_due},
        \texttt{routine\_completed}, \texttt{routine\_skipped}, \texttt{evidence\_submitted},
        \texttt{incident\_opened}, \texttt{incident\_step\_completed},
        \texttt{delegation\_active}.
\end{description}

A complete timeline for a pool-cleaning run looks like:

\begin{minted}{yaml}
{"seq":1,"ts":"2025-04-28T08:00:00Z","type":"delegation_active",
 "delegate_name":"Carlos Méndez","routine_id":"pool_clean"}
{"seq":2,"ts":"2025-04-28T08:00:01Z","type":"routine_due",
 "routine_id":"pool_clean","zone":"pool","executor":"Carlos Méndez"}
{"seq":3,"ts":"2025-04-28T09:15:00Z","type":"evidence_submitted",
 "step_id":"pool_clean","evidence_type":"photo",
 "submitted_by":"Carlos Méndez","sha256":"a3f2...c1d9"}
{"seq":4,"ts":"2025-04-28T09:15:01Z","type":"routine_completed",
 "routine_id":"pool_clean","duration_minutes":75}
\end{minted}

% ─────────────────────────────────────────────────────────────────────────────
\section{Evidence Bundles for Compliance}
\label{sec:evidence-bundles}
% ─────────────────────────────────────────────────────────────────────────────

For warranty claims, insurance documentation, or property sale due diligence, evidence can
be exported as a signed bundle. The bundle aggregates all evidence records for a routine or
date range.

\begin{minted}{yaml}
bundle:
  property: Casa Santiago
  generated_at: "2025-05-01T12:00:00Z"
  routine: pool_clean
  date_range:
    from: "2025-01-01"
    to: "2025-05-01"
  runs:
    - date: "2025-04-28"
      executor: Carlos Méndez
      role: delegate
      completed_at: "2025-04-28T09:15:01Z"
      evidence:
        - type: photo
          filename: pool_clean_2025-04-28.jpg
          sha256: a3f2b19e4c7d8e6f1a2b3c4d5e6f7a8b9c0d1e2f3a4b5c6d7e8f9a0b1c2d3e4f5
    - date: "2025-04-21"
      executor: homeowner
      role: homeowner
      completed_at: "2025-04-21T10:30:00Z"
      evidence:
        - type: photo
          filename: pool_clean_2025-04-21.jpg
          sha256: b4c3d2e1f0a9b8c7d6e5f4a3b2c1d0e9f8a7b6c5d4e3f2a1b0c9d8e7f6a5b4c3
\end{minted}

% ─────────────────────────────────────────────────────────────────────────────
\section{SHA-256 Verification}
\label{sec:evidence-sha256}
% ─────────────────────────────────────────────────────────────────────────────

Photo evidence files are SHA-256 hashed at the moment of submission. The hash is stored both
in the \texttt{trace.jsonl} event and in a \texttt{.sha256} sidecar file alongside the image.

To verify a photo has not been tampered with after submission:

\begin{minted}{yaml}
# The .sha256 sidecar contains the stored hash:
# a3f2b19e4c7d8e6f1a2b3c4d5...  pool_clean_2025-04-28.jpg

# Recompute and compare:
sha256sum pool_clean_2025-04-28.jpg
# If the output matches the sidecar, the file is unmodified.
\end{minted}

\begin{tcolorbox}[colback=yellow!10,colframe=orange!70,title=Note]
  SHA-256 verification proves the file has not been modified \emph{after} submission. It does
  not prove the photo was taken at the claimed location or time. For high-stakes compliance
  use cases, pair photo evidence with metadata (\texttt{EXIF} timestamp and GPS coordinates)
  captured at submission time.
\end{tcolorbox}

% ─────────────────────────────────────────────────────────────────────────────
\section{Reading Evidence from a Completed Run}
\label{sec:evidence-reading}
% ─────────────────────────────────────────────────────────────────────────────

To inspect the evidence from the most recent run of any routine, read the
\texttt{trace.jsonl} file under the \texttt{latest} symlink:

\begin{minted}{yaml}
# List all evidence events from the most recent pool_clean run:
# runs/pool_clean/latest/trace.jsonl

# Filter for evidence_submitted events:
{"seq":3,"ts":"2025-04-28T09:15:00Z","type":"evidence_submitted",
 "step_id":"pool_clean","evidence_type":"photo",
 "label":"Pool water and test strip",
 "submitted_by":"Carlos Méndez","role":"delegate",
 "content":{"filename":"pool_clean_2025-04-28.jpg",
             "sha256":"a3f2...c1d9"}}
\end{minted}

The evidence photos are stored at:
\texttt{runs/\{routine-id\}/\{date\}/evidence/\{filename\}}.

Each evidence file is self-contained: the filename, SHA-256 hash, and submitter identity are
all recorded in the trace, so the evidence bundle can be verified and read independently of
the gert runtime.

\chapter{Reference}
\label{ch:reference}

This chapter contains complete field reference tables for all \texttt{gert.home} schema
elements, environment variables, and error codes.

% ─────────────────────────────────────────────────────────────────────────────
\section{Property File Fields}
\label{sec:ref-property-fields}
% ─────────────────────────────────────────────────────────────────────────────

\begin{center}
\small
\begin{tabular}{lllp{5.5cm}}
\textbf{Field} & \textbf{Type} & \textbf{Required} & \textbf{Description} \\
\hline
\texttt{property.name} & string & YES & Display name for the household unit. \\
\texttt{property.zones} & list & NO & List of named zone objects. \\
\texttt{assets} & list & NO & List of tracked physical asset objects. \\
\texttt{routines} & list & NO & List of recurring routine definitions. \\
\texttt{incident\_templates} & list & NO & List of reactive incident workflow templates. \\
\texttt{consumables} & list & NO & List of consumable items on replacement schedules. \\
\texttt{delegations} & list & NO & List of away-mode delegation blocks. \\
\end{tabular}
\end{center}

% ─────────────────────────────────────────────────────────────────────────────
\section{\texttt{home.routine} Field Reference}
\label{sec:ref-routine-fields}
% ─────────────────────────────────────────────────────────────────────────────

\begin{center}
\small
\begin{tabular}{lllp{5.5cm}}
\textbf{Field} & \textbf{Type} & \textbf{Required} & \textbf{Description} \\
\hline
\texttt{id} & string & YES & Unique routine identifier. \\
\texttt{label} & string & YES & Short display name. \\
\texttt{description} & string & NO & Longer description of the task. \\
\texttt{zone} & string & COND & Zone ID. One of \texttt{zone} or \texttt{asset} required. \\
\texttt{asset} & string & COND & Asset ID. One of \texttt{zone} or \texttt{asset} required. \\
\texttt{cadence.every} & string & COND & Fixed interval (e.g., \texttt{7d}, \texttt{2w}). \\
\texttt{cadence.seasonal} & map & COND & Season-keyed intervals. Mutually exclusive with \texttt{cadence.every}. \\
\texttt{evidence.type} & string & NO & Evidence type: \texttt{photo}, \texttt{note}, \texttt{checklist}, \texttt{none}. \\
\texttt{evidence.prompt} & string & NO & Instruction shown to the executor. \\
\texttt{executor.role} & string & NO & Advisory executor role hint. Not enforced in \texttt{home/v0}. \\
\texttt{notifications.remind\_hours\_before} & integer & NO & Hours before due date to send reminder. \\
\texttt{notifications.escalate\_hours\_overdue} & integer & NO & Hours past due date before overdue alert. \\
\end{tabular}
\end{center}

% ─────────────────────────────────────────────────────────────────────────────
\section{\texttt{home.incident} Field Reference}
\label{sec:ref-incident-fields}
% ─────────────────────────────────────────────────────────────────────────────

\begin{center}
\small
\begin{tabular}{lllp{5.5cm}}
\textbf{Field} & \textbf{Type} & \textbf{Required} & \textbf{Description} \\
\hline
\texttt{id} & string & YES & Unique incident template identifier. \\
\texttt{label} & string & YES & Short display name. \\
\texttt{description} & string & NO & Longer description. \\
\texttt{zones} & list & NO & Zone IDs this incident may affect. \\
\texttt{assets} & list & NO & Asset IDs this incident may affect. \\
\texttt{steps} & list & YES & List of IncidentStep objects. \\
\end{tabular}
\end{center}

% ─────────────────────────────────────────────────────────────────────────────
\section{IncidentStep Field Reference}
\label{sec:ref-incident-step-fields}
% ─────────────────────────────────────────────────────────────────────────────

\begin{center}
\small
\begin{tabular}{lllp{5.5cm}}
\textbf{Field} & \textbf{Type} & \textbf{Required} & \textbf{Description} \\
\hline
\texttt{id} & string & YES & Unique step identifier within the incident. \\
\texttt{type} & string & YES & Step type: \texttt{human\_task}, \texttt{decision}, \texttt{parallel}. \\
\texttt{label} & string & YES & Short display name. \\
\texttt{description} & string & NO & Longer description. \\
\texttt{evidence} & object & NO & Evidence requirement (type + prompt). \\
\texttt{depends\_on} & list & NO & IDs of steps that must complete before this step starts. \\
\texttt{choices} & list & COND & Required for \texttt{type: decision}. List of choice objects. \\
\texttt{parallel\_steps} & list & COND & Required for \texttt{type: parallel}. Sub-steps executed concurrently. \\
\end{tabular}
\end{center}

% ─────────────────────────────────────────────────────────────────────────────
\section{\texttt{home.delegate} Field Reference}
\label{sec:ref-delegate-fields}
% ─────────────────────────────────────────────────────────────────────────────

\begin{center}
\small
\begin{tabular}{lllp{5.5cm}}
\textbf{Field} & \textbf{Type} & \textbf{Required} & \textbf{Description} \\
\hline
\texttt{delegate.name} & string & YES & Display name of the delegate. \\
\texttt{delegate.contact} & string & YES & Phone number or primary contact method. \\
\texttt{delegate.email} & string & NO & Email address of the delegate. \\
\texttt{active.from} & string & YES & Delegation start date (\texttt{YYYY-MM-DD}). \\
\texttt{active.to} & string & YES & Delegation end date (\texttt{YYYY-MM-DD}). \\
\texttt{assigns} & list & YES & Scoped assignments (\texttt{\{routine: id\}} or \texttt{\{zone: id\}}). \\
\texttt{permissions.can\_report\_incidents} & boolean & NO & Whether delegate may open incidents. Default: \texttt{false}. \\
\texttt{permissions.can\_modify\_routines} & boolean & NO & Whether delegate may edit routines. Default: \texttt{false}. \\
\texttt{permissions.can\_view\_history} & boolean & NO & Whether delegate may read run history. Default: \texttt{false}. \\
\texttt{notifications.remind\_delegate\_hours\_before} & integer & NO & Hours before due date to remind delegate. \\
\texttt{notifications.notify\_owner\_on\_completion} & boolean & NO & Notify homeowner when delegate completes a task. \\
\texttt{notifications.notify\_owner\_if\_overdue} & boolean & NO & Notify homeowner if delegate's task is overdue. \\
\end{tabular}
\end{center}

% ─────────────────────────────────────────────────────────────────────────────
\section{Evidence Type Field Reference}
\label{sec:ref-evidence-fields}
% ─────────────────────────────────────────────────────────────────────────────

\subsection{Common Fields}

\begin{description}
  \item[\texttt{type}] Required. One of: \texttt{photo}, \texttt{note},
        \texttt{checklist}, \texttt{none}.
  \item[\texttt{prompt}] Optional. Instruction or question shown to the executor at the
        time of evidence submission.
\end{description}

\subsection{\texttt{photo}}

Requires submission of an image file. The file is SHA-256 hashed and stored in the run
directory. Use for outdoor and physical tasks where completion is visually verifiable.

\subsection{\texttt{note}}

Requires submission of a free-text string. Use for recording cost, measurements, or
contractor communications.

\subsection{\texttt{checklist}}

Requires completion of a structured list of items. Each item is checked individually before
the step closes. Use for multi-step verification tasks.

\subsection{\texttt{none}}

No evidence required. The step closes on executor confirmation. Use sparingly --- only when
no record is needed and completion is self-evident.

% ─────────────────────────────────────────────────────────────────────────────
\section{Consumable Field Reference}
\label{sec:ref-consumable-fields}
% ─────────────────────────────────────────────────────────────────────────────

\begin{center}
\small
\begin{tabular}{lllp{5.5cm}}
\textbf{Field} & \textbf{Type} & \textbf{Required} & \textbf{Description} \\
\hline
\texttt{id} & string & YES & Unique consumable identifier. \\
\texttt{label} & string & YES & Display name. \\
\texttt{description} & string & NO & Longer description. \\
\texttt{zone} & string & NO & Zone where the consumable is used. \\
\texttt{asset} & string & NO & Asset ID this consumable belongs to. \\
\texttt{replace\_every} & string & YES & Replacement interval (e.g., \texttt{180d}). \\
\texttt{last\_replaced} & string & YES & Date of last replacement (\texttt{YYYY-MM-DD}). \\
\texttt{stock.current} & integer & NO & Current stock quantity on hand. \\
\texttt{stock.reorder\_at} & integer & NO & Stock level that triggers a reorder alert. \\
\texttt{triggers\_routine} & string & NO & Routine ID to fire when replacement is due. \\
\end{tabular}
\end{center}

% ─────────────────────────────────────────────────────────────────────────────
\section{Environment Variables}
\label{sec:ref-env-vars}
% ─────────────────────────────────────────────────────────────────────────────

\begin{description}
  \item[\texttt{GERT\_HOME\_PROPERTY\_FILE}] Path to the property file. Default:
        \texttt{./casa.home.yaml}.
  \item[\texttt{GERT\_HOME\_RUN\_DIR}] Root directory for run output and evidence storage.
        Default: \texttt{./runs}.
  \item[\texttt{GERT\_HOME\_TZ}] Timezone used for cadence computation and delegation
        date range evaluation. Default: \texttt{UTC}. Use IANA timezone names
        (e.g., \texttt{America/Santiago}).
\end{description}

% ─────────────────────────────────────────────────────────────────────────────
\section{Error Codes}
\label{sec:ref-error-codes}
% ─────────────────────────────────────────────────────────────────────────────

\begin{center}
\small
\begin{tabular}{llp{7cm}}
\textbf{Code} & \textbf{Stage} & \textbf{Description} \\
\hline
\texttt{HOME\_ERR\_001} & Parse & Property file not found at the path specified by
                                  \texttt{GERT\_HOME\_PROPERTY\_FILE}. \\
\texttt{HOME\_ERR\_002} & Validate & Invalid cadence format. The \texttt{every} value must
                                     be a positive integer followed by \texttt{d} or \texttt{w}
                                     (e.g., \texttt{7d}, \texttt{2w}). \\
\texttt{HOME\_ERR\_003} & Validate & Zone reference not found. A routine, incident template,
                                     or delegation references a zone ID that is not declared
                                     in \texttt{property.zones}. \\
\texttt{HOME\_ERR\_004} & Validate & Asset reference not found. A routine or consumable
                                     references an asset ID that is not declared in
                                     \texttt{assets}. \\
\texttt{HOME\_ERR\_005} & Parse & Delegation date parse error. The \texttt{active.from}
                                  or \texttt{active.to} value is not a valid \texttt{YYYY-MM-DD}
                                  date string. \\
\end{tabular}
\end{center}


\end{document}
